\documentclass[12pt]{article}
\usepackage{amsmath, amsthm, amssymb}
\usepackage{geometry}
\geometry{letterpaper, margin=1in}

\newtheorem{theorem}{Theorem}
\newtheorem{lemma}{Lemma}
\newtheorem{definition}{Definition}

\title{Band-Limitedness of $Z(t)$ in the $\theta$-Variable}
\author{Stephen Crowley}
\date{April 19, 2026}

\begin{document}

\maketitle

\begin{abstract}
Let $Z:\mathbb{R}\to\mathbb{R}$ denote Hardy's function, defined by $\zeta(\tfrac{1}{2}+it) = e^{-i\theta(t)}Z(t)$, where $\theta(t) = \Im\log\Gamma\!\left(\tfrac{1}{4}+\tfrac{it}{2}\right) - \tfrac{t}{2}\log\pi$ is the Riemann--Siegel theta function. The truncated transform $K_T(\nu) := \frac{1}{2\pi}\int_{0}^{T}\zeta(\tfrac{1}{2}+it)\sqrt{\theta'(t)}\,e^{-i\nu\theta(t)}\,dt$ is shown to satisfy $\lim_{T\to\infty}K_T(\nu)=0$ pointwise for every $\nu\in\mathbb{R}\setminus[-2,0]$. The proof rests on exact algebraic cancellation of $\theta''(t(x))$ in the integration-by-parts quotient under the substitution $x=\theta'(t)$, reducing each Riemann--Siegel mode integral to a rational kernel whose decay is controlled to $O(\theta'(T)^{-1/2})$.
\end{abstract}

\section{Introduction}

Throughout this paper, $\zeta$ denotes the Riemann zeta function and $Z$ denotes Hardy's function, related by $\zeta(\tfrac{1}{2}+it)=e^{-i\theta(t)}Z(t)$ with $\theta(t)=\Im\log\Gamma(\tfrac{1}{4}+\tfrac{it}{2})-\tfrac{t}{2}\log\pi$. Fix $T_0>0$ such that $\theta'(T_0)>0$ and $N(T_0)\ge 1$, where $N(t):=\lfloor\sqrt{t/(2\pi)}\rfloor$.

\begin{definition}
For $T\ge T_0$ and $\nu\in\mathbb{R}$, define the truncated $\theta$-transform of $Z$ by
\begin{equation}
K_T(\nu) := \frac{1}{2\pi}\int_{T_0}^{T} Z(t)\sqrt{\theta'(t)}\,e^{-i(\nu+1)\theta(t)}\,dt.
\end{equation}
\end{definition}

The substitution $\zeta(\tfrac{1}{2}+it)=e^{-i\theta(t)}Z(t)$ identifies this with $\frac{1}{2\pi}\int_{T_0}^{T}\zeta(\tfrac{1}{2}+it)\sqrt{\theta'(t)}\,e^{-i\nu\theta(t)}\,dt$. Set $\mu:=\nu+1$, so that the condition $\nu\notin[-2,0]$ is equivalent to $|\mu|>1$.

The principal result stated below establishes pointwise vanishing of $K_T(\nu)$ as $T\to\infty$ for every $\nu\in\mathbb{R}\setminus[-2,0]$.

\section{Statement and Proof}

\begin{theorem}
For every fixed $\nu\in\mathbb{R}\setminus[-2,0]$,
\begin{equation}
\lim_{T\to\infty} K_T(\nu) = 0.
\end{equation}
\end{theorem}

\begin{proof}
The digamma series $\psi'(s)=\sum_{k\ge 0}(s+k)^{-2}$ evaluated at $s=\tfrac{1}{4}+\tfrac{it}{2}$ yields
\begin{equation}
\theta''(t) = \frac{1}{4}\sum_{k\ge 0}\frac{2(\tfrac{1}{4}+k)\cdot\tfrac{t}{2}}{\bigl|\tfrac{1}{4}+k+\tfrac{it}{2}\bigr|^{4}} > 0 \qquad (t>0).
\end{equation}
Each summand is strictly positive for $t>0$, hence $\theta':[T_0,\infty)\to[\theta'(T_0),\infty)$ is a strictly increasing $C^{\infty}$ bijection with $C^{\infty}$ inverse $t:[\theta'(T_0),\infty)\to[T_0,\infty)$, $x\mapsto t(x)$. Moreover $\theta'(t)\to\infty$ as $t\to\infty$ via $\theta'(t)=\tfrac{1}{2}\log(t/(2\pi))+O(t^{-2})$.

The Riemann--Siegel identity provides the exact decomposition $Z(t)=2\sum_{n=1}^{N(t)}n^{-1/2}\cos(\theta(t)-t\log n)+R(t)$, where $R(t)=O(t^{-1/4})$ is the Riemann--Siegel remainder. Substituting $\cos\alpha=\tfrac{1}{2}(e^{i\alpha}+e^{-i\alpha})$ and exchanging the finite sum with the integral yields
\begin{equation}
K_T(\nu) = \frac{1}{2\pi}\sum_{\sigma\in\{+1,-1\}}\sum_{n\le N(T)} n^{-1/2}\, J_{n,\sigma}(T,\mu) + \mathcal{R}(T,\mu),
\end{equation}
where, with $t_n:=\max(T_0,2\pi n^{2})$,
\begin{equation}
J_{n,\sigma}(T,\mu) := \int_{t_n}^{T}\sqrt{\theta'(t)}\,e^{i[(\sigma-\mu)\theta(t)-\sigma t\log n]}\,dt,
\end{equation}
and $\mathcal{R}(T,\mu):=\frac{1}{2\pi}\int_{T_0}^{T}R(t)\sqrt{\theta'(t)}\,e^{-i\mu\theta(t)}\,dt$. Since $R(t)\sqrt{\theta'(t)}=O(t^{-1/4}(\log t)^{1/2})$, integration by parts in $\theta(t)$ against the nonstationary phase $e^{-i\mu\theta(t)}$ (the phase derivative is $-\mu\theta'(t)\to-\infty$, hence bounded away from zero on $[T_0,\infty)$) shows $\mathcal{R}(T,\mu)$ converges absolutely as $T\to\infty$ and the corresponding tail $\mathcal{R}(\infty,\mu)-\mathcal{R}(T,\mu)$ is $O(T^{-1/4}(\log T)^{-1/2})\to 0$.

Set $X:=\theta'(T)$, $\beta_n:=\theta'(2\pi n^{2})$, and $x_0:=\theta'(T_0)$. The substitution $x=\theta'(t)$, $dx=\theta''(t)\,dt$, transforms each mode integral to
\begin{equation}
J_{n,\sigma}(T,\mu) = \int_{\beta_n\vee x_0}^{X}\frac{\sqrt{x}}{\theta''(t(x))}\,e^{i\widetilde{\Phi}_{n,\sigma,\mu}(x)}\,dx,
\end{equation}
where the phase $\widetilde{\Phi}_{n,\sigma,\mu}(x):=(\sigma-\mu)\theta(t(x))-\sigma t(x)\log n$ has derivative
\begin{equation}
\widetilde{\Phi}'_{n,\sigma,\mu}(x) = \frac{(\sigma-\mu)x - \sigma\log n}{\theta''(t(x))}.
\end{equation}
The numerator is affine in $x$ with unique zero $x^{*}_{n,\sigma}=\sigma\log n/(\sigma-\mu)$. The amplitude-to-phase-derivative ratio admits exact cancellation of $\theta''(t(x))$:
\begin{equation}
\frac{A_{n}(x)}{\widetilde{\Phi}'_{n,\sigma,\mu}(x)} = \frac{\sqrt{x}/\theta''(t(x))}{[(\sigma-\mu)x-\sigma\log n]/\theta''(t(x))} = \frac{\sqrt{x}}{(\sigma-\mu)x - \sigma\log n} =: Q_{n,\sigma,\mu}(x),
\end{equation}
a function of $x$ alone. Its derivative is
\begin{equation}
Q'_{n,\sigma,\mu}(x) = -\frac{(\sigma-\mu)x + \sigma\log n}{2\sqrt{x}\,[(\sigma-\mu)x - \sigma\log n]^{2}}.
\end{equation}

For $|\mu|>1$, the unique sign-condition zero $x^{*}_{n,\sigma}>0$ occurs in exactly one of the two branches $\sigma\in\{+1,-1\}$, and equals $\log n/(1+|\mu|)$. Define
\begin{equation}
S(\mu) := \bigl\{n\in\mathbb{N}\,:\,\beta_n\le\tfrac{\log n}{1+|\mu|}\bigr\}.
\end{equation}

\begin{lemma}
$S(\mu)$ has finite cardinality for every $|\mu|>1$.
\end{lemma}
\begin{proof}
The asymptotic $\psi(z)\sim\log z$ implies $\theta'(t)=\tfrac{1}{2}\log(t/(2\pi))+O(t^{-2})$, hence $\beta_n=\theta'(2\pi n^{2})=\log n + O(n^{-4})$. Therefore
\begin{equation}
\beta_n - \tfrac{\log n}{1+|\mu|} = \tfrac{|\mu|}{1+|\mu|}\log n + O(n^{-4}) \longrightarrow +\infty.
\end{equation}
There exists $N_{0}(\mu)<\infty$ such that $n>N_{0}(\mu)\Rightarrow n\notin S(\mu)$, so $|S(\mu)|\le N_{0}(\mu)<\infty$.
\end{proof}

For $n\notin S(\mu)$, the zero $x^{*}_{n,\sigma}$ lies strictly below $\beta_n\vee x_0$, hence $\widetilde{\Phi}'_{n,\sigma,\mu}$ has constant sign on $[\beta_n\vee x_0,X]$ for every $X>\beta_n\vee x_0$, and the denominator $(\sigma-\mu)x-\sigma\log n$ is bounded away from zero. Integration by parts with $u=Q_{n,\sigma,\mu}(x)$, $dv=i\widetilde{\Phi}'_{n,\sigma,\mu}(x)\,e^{i\widetilde{\Phi}_{n,\sigma,\mu}(x)}\,dx$ yields
\begin{align}
J_{n,\sigma}(T,\mu) =\;& \frac{\sqrt{X}}{i\,[(\sigma-\mu)X-\sigma\log n]}\,e^{i\widetilde{\Phi}_{n,\sigma,\mu}(X)} \notag\\
&-\frac{\sqrt{\beta_n\vee x_0}}{i\,[(\sigma-\mu)(\beta_n\vee x_0)-\sigma\log n]}\,e^{i\widetilde{\Phi}_{n,\sigma,\mu}(\beta_n\vee x_0)} \notag\\
&+\int_{\beta_n\vee x_0}^{X}\frac{(\sigma-\mu)x+\sigma\log n}{2i\sqrt{x}\,[(\sigma-\mu)x-\sigma\log n]^{2}}\,e^{i\widetilde{\Phi}_{n,\sigma,\mu}(x)}\,dx,
\end{align}
which is the exact identity. The boundary at $X$ satisfies
\begin{equation}
\left|\frac{\sqrt{X}}{(\sigma-\mu)X-\sigma\log n}\right| \le \frac{\sqrt{X}}{(|\mu|-1)X - |\log n|} \;\sim\; \frac{1}{(|\mu|-1)\sqrt{X}} \longrightarrow 0
\end{equation}
as $X\to\infty$. The integrand of the remainder satisfies the absolute bound
\begin{equation}
\left|\frac{(\sigma-\mu)x+\sigma\log n}{2\sqrt{x}\,[(\sigma-\mu)x-\sigma\log n]^{2}}\right| = O(x^{-3/2}) \qquad (x\to\infty),
\end{equation}
so the remainder integral converges absolutely on $[\beta_n\vee x_0,\infty)$ and the tail $\int_{X}^{\infty}|\cdot|\,dx = O(X^{-1/2})\to 0$.

For $n\in S(\mu)$, the stationary point $t^{*}_{n,\sigma}=t(x^{*}_{n,\sigma})$ is a fixed finite value determined by $\theta'(t^{*}_{n,\sigma})=\log n/(1+|\mu|)$. The mode $J_{n,\sigma}(T,\mu)$ converges as $T\to\infty$ to a finite limit $J_{n,\sigma}(\infty,\mu)\in\mathbb{C}$, contributing to $K_{\infty}(\nu)$. The quantity to control is the tail
\begin{equation}
J_{n,\sigma}(\infty,\mu)-J_{n,\sigma}(T,\mu) = \int_{T}^{\infty}\sqrt{\theta'(t)}\,e^{i\Phi_{n,\sigma,\mu}(t)}\,dt,
\end{equation}
where $\Phi_{n,\sigma,\mu}(t)=(\sigma-\mu)\theta(t)-\sigma t\log n$. For $T>t^{*}_{n,\sigma}$, strict monotonicity of $\theta'$ gives $\theta'(T)>\log n/(1+|\mu|)$, whence
\begin{equation}
|\Phi'_{n,\sigma,\mu}(t)| = |(\sigma-\mu)\theta'(t)-\sigma\log n| \ge (|\mu|-1)\bigl(\theta'(T)-\tfrac{\log n}{|\mu|-1}\bigr) > 0
\end{equation}
for all $t\ge T$, with the lower bound diverging as $T\to\infty$. The denominator in the integration-by-parts quotient is therefore bounded away from zero on $[T,\infty)$, and the argument applied to $n\notin S(\mu)$ transfers verbatim to the tail integral, yielding upper boundary $O(\theta'(T)^{-1/2})$ and remainder $O(\theta'(T)^{-1/2})$, both vanishing as $T\to\infty$.

The cardinality $|S(\mu)|$ is finite, so the aggregate contribution from $\{n\in S(\mu)\}$ is a finite sum of terms each tending to zero. Combined with the vanishing of $\mathcal{R}(\infty,\mu)-\mathcal{R}(T,\mu)$ and of the $n\notin S(\mu)$ tails, every component of $K_{\infty}(\nu)-K_T(\nu)$ tends to zero as $T\to\infty$. Therefore
\begin{equation}
\lim_{T\to\infty} K_T(\nu) = 0 \qquad \text{for every fixed } \nu\in\mathbb{R}\setminus[-2,0].
\end{equation}
\end{proof}

\end{document}
